\documentclass[10pt,a4paper]{amsart}
\usepackage{amsmath,amssymb,amsthm,graphicx}
% Retained minimal preamble for the Stacks Project chapter flat.tex (More on Flatness).
% Verbatim excerpt of the chapter preamble: Xy-pic setup (native preamble lines 14-18) and the native
% theorem declarations with the mathematical macros (native preamble lines 154-217). Copied mechanically;
% nothing added, removed or redefined. The wrapper class and the amsmath/amssymb/amsthm/graphicx base
% packages are supplied by the assembler, so the class/externaldocument part of the native preamble is omitted.
% For commutative diagrams we use Xy-pic
\usepackage[all]{xy}

% We use 2cell for 2-commutative diagrams.
\xyoption{2cell}
% Theorem environments.
%
\theoremstyle{plain}
\newtheorem{theorem}[subsection]{Theorem}
\newtheorem{proposition}[subsection]{Proposition}
\newtheorem{lemma}[subsection]{Lemma}

\theoremstyle{definition}
\newtheorem{definition}[subsection]{Definition}
\newtheorem{example}[subsection]{Example}
\newtheorem{exercise}[subsection]{Exercise}
\newtheorem{situation}[subsection]{Situation}

\theoremstyle{remark}
\newtheorem{remark}[subsection]{Remark}
\newtheorem{remarks}[subsection]{Remarks}

\numberwithin{equation}{subsection}

% Macros
%
\def\lim{\mathop{\mathrm{lim}}\nolimits}
\def\colim{\mathop{\mathrm{colim}}\nolimits}
\def\Spec{\mathop{\mathrm{Spec}}}
\def\Hom{\mathop{\mathrm{Hom}}\nolimits}
\def\Ext{\mathop{\mathrm{Ext}}\nolimits}
\def\SheafHom{\mathop{\mathcal{H}\!\mathit{om}}\nolimits}
\def\SheafExt{\mathop{\mathcal{E}\!\mathit{xt}}\nolimits}
\def\Sch{\mathit{Sch}}
\def\Mor{\mathop{\mathrm{Mor}}\nolimits}
\def\Ob{\mathop{\mathrm{Ob}}\nolimits}
\def\Sh{\mathop{\mathit{Sh}}\nolimits}
\def\NL{\mathop{N\!L}\nolimits}
\def\CH{\mathop{\mathrm{CH}}\nolimits}
\def\proetale{{pro\text{-}\acute{e}tale}}
\def\etale{{\acute{e}tale}}
\def\QCoh{\mathit{QCoh}}
\def\Ker{\mathop{\mathrm{Ker}}}
\def\Im{\mathop{\mathrm{Im}}}
\def\Coker{\mathop{\mathrm{Coker}}}
\def\Coim{\mathop{\mathrm{Coim}}}

% Boxtimes
%
\DeclareMathSymbol{\boxtimes}{\mathbin}{AMSa}{"02}

%
% Macros for moduli stacks/spaces
%
\def\QCohstack{\mathcal{QC}\!\mathit{oh}}
\def\Cohstack{\mathcal{C}\!\mathit{oh}}
\def\Spacesstack{\mathcal{S}\!\mathit{paces}}
\def\Quotfunctor{\mathrm{Quot}}
\def\Hilbfunctor{\mathrm{Hilb}}
\def\Curvesstack{\mathcal{C}\!\mathit{urves}}
\def\Polarizedstack{\mathcal{P}\!\mathit{olarized}}
\def\Complexesstack{\mathcal{C}\!\mathit{omplexes}}
% \Pic is the operator that assigns to X its picard group, usage \Pic(X)
% \Picardstack_{X/B} denotes the Picard stack of X over B
% \Picardfunctor_{X/B} denotes the Picard functor of X over B
\def\Pic{\mathop{\mathrm{Pic}}\nolimits}
\def\Picardstack{\mathcal{P}\!\mathit{ic}}
\def\Picardfunctor{\mathrm{Pic}}
\def\Deformationcategory{\mathcal{D}\!\mathit{ef}}

% --- appended 2026-10-12 (Chief): the project own definitions of the reference/slogan/history
% display environments, copied mechanically from the shared Stacks preamble lines 8-12
% (\usepackage{verbatim} plus the three \newenvironment lines, which the chapter preamble keeps
% verbatim from the project preamble). Nothing added, removed or redefined: these environments
% carry no mathematical content and the project itself typesets them as comments. Without these
% three lines a counted statement whose own text uses them fails with Environment slogan undefined.
\usepackage{verbatim}
\newenvironment{reference}{\comment}{\endcomment}
\newenvironment{slogan}{\comment}{\endcomment}
\newenvironment{history}{\comment}{\endcomment}
\usepackage{cleveref}
\crefname{originallemma}{original lemma}{original lemmas}
\Crefname{originallemma}{Original lemma}{Original lemmas}
\makeatletter
\newlabel{algebra-lemma-ass-quotient-ring}{{05BY}{source}{}{sourcepointer-1}{}}
\newlabel{algebra-lemma-ass-quotient-ring@cref}{{[originallemma][0][]05BY}{[1][0][]source}{}{}{}}
\makeatother
\crefname{originallemma}{original lemma}{original lemmas}
\Crefname{originallemma}{Original lemma}{Original lemmas}
\makeatletter
\newlabel{divisors-lemma-weakly-associated-finite}{{05EZ}{source}{}{sourcepointer-2}{}}
\newlabel{divisors-lemma-weakly-associated-finite@cref}{{[originallemma][0][]05EZ}{[1][0][]source}{}{}{}}
\makeatother
\crefname{originallemma}{original lemma}{original lemmas}
\Crefname{originallemma}{Original lemma}{Original lemmas}
\makeatletter
\newlabel{lemma-etale-weak-assassin-up-down}{{05FP}{source}{}{sourcepointer-3}{}}
\newlabel{lemma-etale-weak-assassin-up-down@cref}{{[originallemma][0][]05FP}{[1][0][]source}{}{}{}}
\makeatother
\crefname{originallemma}{original lemma}{original lemmas}
\Crefname{originallemma}{Original lemma}{Original lemmas}
\makeatletter
\newlabel{lemma-weak-bourbaki-pre}{{05IG}{source}{}{sourcepointer-4}{}}
\newlabel{lemma-weak-bourbaki-pre@cref}{{[originallemma][0][]05IG}{[1][0][]source}{}{}{}}
\makeatother
\crefname{originallemma}{original lemma}{original lemmas}
\Crefname{originallemma}{Original lemma}{Original lemmas}
\makeatletter
\newlabel{proposition-finite-type-flat-at-point}{{05I5}{source}{}{sourcepointer-5}{}}
\newlabel{proposition-finite-type-flat-at-point@cref}{{[originallemma][0][]05I5}{[1][0][]source}{}{}{}}
\makeatother
\title{Weakly associated points of a tensor product with a pulled back sheaf in the finite type case}
\author{The Stacks Project Authors}
\date{stacks-project:flat.tex:local-snapshot-20261011; GNU FDL 1.2 or later}
\begin{document}
\maketitle
\noindent\textit{Editorial note.} All delivered mathematical text is the author's own, copied from the pinned Stacks Project chapter file flat.tex; no statement or proof step was rewritten, reordered or removed.\par\medskip
\noindent\textit{Editorial note.} Preamble provenance: the chapter's own retained minimal preamble (sources/stacks-flat/retained-preamble.tex, lines 1-85) is a verbatim excerpt of the native Xy-pic setup and of the native theorem declarations and mathematical macros, followed by the project's own reference/slogan/history display-environment definitions copied mechanically from the shared preamble; the wrapper class and the amsmath/amssymb/amsthm/graphicx base packages are supplied by the assembler, and extra\_packages is left empty because everything needed is already in the retained preamble. No macro was invented or redefined.\par\medskip
\noindent\textit{Editorial note.} Same-chapter pointers are located at their \textbackslash{}label lines inside the pinned chapter file; cross-chapter pointers are located in the project's own tag list staged in the snapshot (sources/stacks-chow/tags.txt), at the verbatim lines carrying their canonical tags.\par\medskip
\noindent\textit{Allowed standard prerequisites (editorial).}\begin{itemize}
\item Weakly associated points of a closed subscheme and of a quotient ring. Divisors, label divisors-lemma-weakly-associated-finite (proof line 4017); Algebra, label algebra-lemma-ass-quotient-ring (proof line 4020). Transfers the weakly associated point and the associated point of the fibre back and forth between X and the affine space in which it is embedded, which is the reduction step of the first paragraph.
\item The local structure of a finite type sheaf flat at a point. Same chapter flat.tex, label proposition-finite-type-flat-at-point, invoked at proof line 4029. Produces the elementary etale neighbourhood and the open V over which the pullback of F becomes of finite presentation and flat, which is the reduction of the second paragraph. The proposition is claimed separately as S14-algebra-finite-type-flat-at-point.
\item Weakly associated primes under an etale neighbourhood, and the finite presentation case. Same chapter flat.tex, labels lemma-etale-weak-assassin-up-down (proof line 4033) and lemma-weak-bourbaki-pre (proof line 4056). The first lets the statement be checked after the elementary etale neighbourhood, the second is the already proven finite presentation statement that finishes the proof. The latter is claimed separately as S14-algebra-weak-bourbaki-pre.
\end{itemize}\medskip
\ifdefined\hypertarget\hypertarget{sourcepointer-1}{}\fi
\noindent\textit{Source reference.} Original lemma 05BY: sources/stacks-chow/tags.txt, source lines 6901--6901. Navigation pointer for \textbackslash{}ref\{algebra-lemma-ass-quotient-ring\} used inside the counted span (line 4020); the argument that consumes it is retained in the span.\par\medskip
\ifdefined\hypertarget\hypertarget{sourcepointer-2}{}\fi
\noindent\textit{Source reference.} Original lemma 05EZ: sources/stacks-chow/tags.txt, source lines 7008--7008. Navigation pointer for \textbackslash{}ref\{divisors-lemma-weakly-associated-finite\} used inside the counted span (line 4017); the argument that consumes it is retained in the span.\par\medskip
\ifdefined\hypertarget\hypertarget{sourcepointer-3}{}\fi
\noindent\textit{Source reference.} Original lemma 05FP: sources/stacks-flat/pinned-original-source.tex, source lines 353--353. Navigation pointer for \textbackslash{}ref\{lemma-etale-weak-assassin-up-down\} used inside the counted span (line 4033); the argument that consumes it is retained in the span.\par\medskip
\ifdefined\hypertarget\hypertarget{sourcepointer-4}{}\fi
\noindent\textit{Source reference.} Original lemma 05IG: sources/stacks-flat/pinned-original-source.tex, source lines 3750--3750. Navigation pointer for \textbackslash{}ref\{lemma-weak-bourbaki-pre\} used inside the counted span (line 4056); the argument that consumes it is retained in the span.\par\medskip
\ifdefined\hypertarget\hypertarget{sourcepointer-5}{}\fi
\noindent\textit{Source reference.} Original lemma 05I5: sources/stacks-flat/pinned-original-source.tex, source lines 2477--2477. Navigation pointer for \textbackslash{}ref\{proposition-finite-type-flat-at-point\} used inside the counted span (line 4029); the argument that consumes it is retained in the span.\par\medskip
\begin{lemma}
\label{lemma-bourbaki-finite-type-general-base-at-point}
Let $S$ be a scheme.
Let $f : X \to S$ be locally of finite type.
Let $x \in X$ with image $s \in S$.
Let $\mathcal{F}$ be a finite type quasi-coherent sheaf on $X$.
Let $\mathcal{G}$ be a quasi-coherent sheaf on $S$.
If $\mathcal{F}$ is flat at $x$ over $S$, then
$$
x \in \text{WeakAss}_X(\mathcal{F} \otimes_{\mathcal{O}_X} f^*\mathcal{G})
\Leftrightarrow
s \in \text{WeakAss}_S(\mathcal{G})
\text{ and }
x \in \text{Ass}_{X_s}(\mathcal{F}_s).
$$
\end{lemma}

\begin{proof}
In this paragraph we reduce to $f$ being of finite presentation.
The question is local on $X$ and $S$, hence we may assume $X$ and $S$
are affine. Write $X = \Spec(B)$, $S = \Spec(A)$ and write
$B = A[x_1, \ldots, x_n]/I$. In other words we obtain a closed immersion
$i : X \to \mathbf{A}^n_S$ over $S$. Denote $t = i(x) \in \mathbf{A}^n_S$.
Note that $i_*\mathcal{F}$ is a finite type quasi-coherent sheaf on
$\mathbf{A}^n_S$ which is flat at $t$ over $S$ and note that
$$
i_*(\mathcal{F} \otimes_{\mathcal{O}_X} f^*\mathcal{G}) =
i_*\mathcal{F} \otimes_{\mathcal{O}_{\mathbf{A}^n_S}} p^*\mathcal{G}
$$
where $p : \mathbf{A}^n_S \to S$ is the projection. Note that
$t$ is a weakly associated point of
$i_*(\mathcal{F} \otimes_{\mathcal{O}_X} f^*\mathcal{G})$
if and only if $x$ is a weakly associated point of
$\mathcal{F} \otimes_{\mathcal{O}_X} f^*\mathcal{G}$, see
Divisors, Lemma \ref{divisors-lemma-weakly-associated-finite}.
Similarly $x \in \text{Ass}_{X_s}(\mathcal{F}_s)$ if and only
if $t \in \text{Ass}_{\mathbf{A}^n_s}((i_*\mathcal{F})_s)$ (see
Algebra, Lemma \ref{algebra-lemma-ass-quotient-ring}).
Hence it suffices to prove the lemma in case $X = \mathbf{A}^n_S$.
Thus we may assume that $X \to S$ is of finite presentation.

\medskip\noindent
In this paragraph we reduce to $\mathcal{F}$ being of finite presentation
and flat over $S$.
Choose an elementary \'etale neighbourhood $e : (S', s') \to (S, s)$
and an open $V \subset X \times_S \Spec(\mathcal{O}_{S', s'})$
as in Proposition \ref{proposition-finite-type-flat-at-point}.
Let $x' \in X' = X \times_S S'$ be the unique point mapping to $x$
and $s'$. Then it suffices to prove the statement for
$X' \to S'$, $x'$, $s'$, $(X' \to X)^*\mathcal{F}$, and $e^*\mathcal{G}$, see
Lemma \ref{lemma-etale-weak-assassin-up-down}.
Let $v \in V$ the unique point mapping to $x'$
and let $s' \in \Spec(\mathcal{O}_{S', s'})$ be the closed point.
Then $\mathcal{O}_{V, v} = \mathcal{O}_{X', x'}$
and $\mathcal{O}_{\Spec(\mathcal{O}_{S', s'}), s'} =
\mathcal{O}_{S', s'}$ and similarly for the stalks of pullbacks of
$\mathcal{F}$ and $\mathcal{G}$.
Also $V_{s'} \subset X'_{s'}$ is an open subscheme.
Since the condition of being a weakly associated point
depend only on the stalk of the sheaf, we may
replace
$X' \to S'$, $x'$, $s'$, $(X' \to X)^*\mathcal{F}$, and $e^*\mathcal{G}$
by
$V \to \Spec(\mathcal{O}_{S', s'})$, $v$, $s'$, $(V \to X)^*\mathcal{F}$,
and $(\Spec(\mathcal{O}_{S', s'}) \to S)^*\mathcal{G}$.
Thus we may assume that $f$ is of finite presentation and
$\mathcal{F}$ of finite presentation and flat over $S$.

\medskip\noindent
Assume $f$ is of finite presentation and
$\mathcal{F}$ of finite presentation and flat over $S$.
After shrinking $X$ and $S$ to affine neighbourhoods
of $x$ and $s$, this case is handled by
Lemma \ref{lemma-weak-bourbaki-pre}.
\end{proof}

\end{document}
